\section{Reconciling the two weight-seven Gaussian candidates}
\label{s6audit:sec:main}

The pinned manuscript's Conjecture~\texttt{cycloquot:conj:S6} and the
incoming Cayley continuation state different-looking formulas for $S_6$.
They are exactly equivalent. This fact neither proves nor refutes the
common conjecture, but it prevents the same proposed identity from being
counted twice and removes an unnecessary product from one presentation.
Write $G=\beta(2)$ and $L=\log2$, and recall
\[
 S_6=\sum_{n\ge0}\frac{(-1)^nH_n}{(2n+1)^6},
 \qquad g_{a,b}=\operatorname{Im}\operatorname{Li}_{a,b}(i,1).
\]
The canonical candidate is
\begin{align}
 S_6\stackrel{?}{=}&\frac{722}{527}g_{6,1}
 +\frac{40}{527}g_{4,3}-\frac{128}{155}g_{2,5}
 +\frac{15191}{28569600}\pi^7\notag\\
 &-\frac3{124}G\zeta(5)-\frac{2373}{10540}\beta(4)\zeta(3)
 -2\beta(6)L.\label{s6audit:eq:old-candidate}
\end{align}
It was supported in the baseline by two independent high-precision
methods after its coefficients were frozen. Those numerical records
are inherited evidence, not a new verification performed here.
The present exact calculation concerns the equivalence of
\eqref{s6audit:eq:old-candidate} and the incoming Cayley presentation
\cite{cont:CayleyIncoming}.

\begin{proposition}[Exact change of Gaussian coordinates]\label{s6audit:prop:coordinate}
With $g_{a,b}=\operatorname{Im}\operatorname{Li}_{a,b}(i,1)$,
\begin{align}
 g_{2,5}={}&30g_{6,1}+15g_{5,2}+5g_{4,3}
 -\frac{11\pi^7}{368640}
 -\frac{15}{512}G\zeta(5)
 +\frac3{16}\beta(4)\zeta(3).
 \label{s6audit:eq:coordinate}
\end{align}
Consequently the manuscript's $S_6$ conjecture is equivalent to
\begin{align}
 S_6\stackrel{?}{=}&-\frac{12334}{527}g_{6,1}
        -\frac{384}{31}g_{5,2}
        -\frac{2136}{527}g_{4,3}
        +\frac{3179}{5713920}\pi^7\notag\\
       &-\frac{801}{2108}\beta(4)\zeta(3)
        -2\beta(6)L.
 \label{s6audit:eq:reduced-candidate}
\end{align}
The latter is precisely the incoming Cayley candidate after dividing its
stored primitive integer vector by its $S_6$ coefficient.
\end{proposition}

\begin{proof}
The convergent integral-shuffle products, valid first for $|z|<1$ and
then by ordinary convergence at $z=i$, are
\begin{align*}
 \operatorname{Li}_2(z)\operatorname{Li}_5(z)
 &=10F_{6,1}(z)+5F_{5,2}(z)+3F_{4,3}(z)
   +2F_{3,4}(z)+F_{2,5}(z),\\
 \operatorname{Li}_3(z)\operatorname{Li}_4(z)
 &=20F_{6,1}(z)+10F_{5,2}(z)+4F_{4,3}(z)+F_{3,4}(z),
\end{align*}
where $F_{a,b}(z)=\operatorname{Li}_{a,b}(z,1)$.
Subtract twice the second equation from the first. Since
\begin{align*}
 \operatorname{Im}(\operatorname{Li}_2(i)\operatorname{Li}_5(i))
 &=-\frac{5\pi^7}{73728}-\frac{15}{512}G\zeta(5),\\
 \operatorname{Im}(\operatorname{Li}_3(i)\operatorname{Li}_4(i))
 &=-\frac{7\pi^7}{368640}-\frac3{32}\beta(4)\zeta(3),
\end{align*}
taking imaginary parts gives \eqref{s6audit:eq:coordinate}.
Substitution in the manuscript candidate cancels its entire
$G\zeta(5)$ coefficient and gives \eqref{s6audit:eq:reduced-candidate}.

For an explicit relation certificate, let $R_{\rm old}$ and
$R_{\rm Cayley}$ denote the two residuals normalized to coefficient
$+1$ on $S_6$. Let $E_{p,q}$ denote the shuffle sum minus
$\operatorname{Im}(\operatorname{Li}_p(i)\operatorname{Li}_q(i))$.
Then, as an exact formal coefficient identity,
\[
 R_{\rm old}-R_{\rm Cayley}
 =\frac{128}{155}\bigl(E_{2,5}-2E_{3,4}\bigr).
\]
The right-hand side vanishes by two convergent shuffles. The
dependency-free program \texttt{verify\_s6\_equivalence.py} rebuilds the
binomial shuffle rows and verifies this coefficient identity and the
direct substitution independently, using exact rational arithmetic.
\end{proof}

\paragraph{Scope of the search.}
A new exploratory system, within the broader standard-relation
setting discussed in \cite{cont:Zhao}, was generated using the pinned incoming
word convention. It contains $51\,244$ nonzero projected weight-seven
double-shuffle rows and three convergent Cayley rows, with initial
imaginary coordinates drawn from $24\,514$ possible convergent words.
The double-shuffle inputs are convergent except for the permitted single
factor $\operatorname{Li}_1(1)$; its common divergence cancels between
shuffle and stuffle. The Cayley seeds are the words of
$F_{6,1}(i),F_{5,2}(i),F_{4,3}(i)$.

Sparse elimination modulo $1\,000\,003$ was stopped for resource control
after $8\,000$ logged pivots, while the target residual remained
nonzero and fill-in had exceeded $20$ million coefficient entries.
This incomplete search is inconclusive: it proves neither membership
nor nonmembership in the row span. It provides no proof of the $S_6$
conjecture, and the word ``proved'' should not be attached to either
equivalent $S_6$ presentation. A future search should exploit shuffle
indecomposables or a better sparse elimination ordering before expanding
the entire weight-seven standard-relation space.
